\subsection{弥散测度；原子测度}

定义 5. --- 称 T 上的一个测度 \(\theta\) 为弥散的，如果 \(|\theta|(\{t\}) = 0\) 对每个 \(t \in \mathrm{T}\) 成立。

例. --- \(\mathbf{R}\) 上的 Lebesgue 测度是弥散的（Ch. IV, §1, No.~3, 注 1）。

说 \(\theta\) 是 T 上的一个弥散测度，等于说每个余集有限的集都承载 \(|\theta|\) ，或者说，\(|\theta|\) 与每个点测度互奇异。因此弥散测度在 \(\mathcal{M}(\mathrm{T})\) 中构成一个带（Ch. II, §1, No.~5, 定理 1）。

回忆（Ch. III, §1, No.~3），称 T 上的一个复测度 \(\rho\) 为原子测度，如果它具有 \(\sum_{t\in \mathrm{T}}\alpha (t)\varepsilon_t\) 的形式，其中 \(\alpha\) 是 T 上的一个复函数，使得 \(\sum_{t\in \mathrm{K}}|\alpha (t)| < +\infty\) 对 T 的每个紧子集 K 成立，这就表达了族 \(\left(\alpha (t)\varepsilon_t\right)_{t\in \mathrm{T}}\) 是可和的（§2, No.~1, 注 2）。于是由 No.~5 的定理 2 的推论 3 之后的评注推出 \(|\rho| = \sum_{t\in \mathrm{T}}|\alpha (t)|\varepsilon_t\) 。出现在这些公式中的函数 \(\alpha\) 是唯一确定的，因为 \(\alpha (t) = \rho (\{t\})\) 。一个原子测度与一个弥散测度互奇异。

命题 15. --- T 上的每个复测度 \(\sigma\) 都可以以一种且仅一种方式写成 \(\rho + \theta\) 的形式，其中 \(\rho\) 是一个原子测度，\(\theta\) 是一个弥散测度；此时有 \(|\sigma| = |\rho| + |\theta|\) 。

分解的唯一性是明显的，因为当 \(\rho\) 是原子的而 \(\theta\) 是弥散的时，必然有 \(\rho = \sum_{t\in \mathrm{T}}\rho (\{t\})\varepsilon_t = \sum_{t\in \mathrm{T}}\sigma (\{t\})\varepsilon_t\) 以及 \(\theta = \sigma -\rho\) 。为证明存在性，只须注意到 \(\sum_{t\in \mathrm{K}}|\sigma (\{t\})| \leqslant |\sigma |(\mathrm{K}) < +\infty\) 对每个紧集 \(\mathbf{K}\) 成立，于是可以置 \(\sum_{t\in \mathrm{T}}\sigma (\{t\})\varepsilon_t = \rho\) 。测度 \(\sigma -\rho\) 显然是弥散的，而关系 \(|\sigma | = |\rho | + |\sigma -\rho |\) 由 No.~7 的命题 13 的推论 3 立即得出。

我们注意到，这个证明表明：若 \(\sigma\) 由一个集 \(M\) 承载，且 \(|\sigma| (\{t\}) > 0\) 对每个 \(t \in M\) 成立，则 \(\sigma\) 是原子的。

